% Burst-aware renewal dynamics --- plan \S4.3.
% Sources: new_notes3 Part D \S18--\S19 (authoritative); new_notes Part II
% \S7--\S8; new_notes2 \S6--\S7.  The universal skeleton of draft \S6.1 is
% promoted in here; the overlays of draft \S3.4 move to \S4.
% Numerics: new_notes3/verify_result_20_1.py (54 checks, all passing).

\section{Burst-aware renewal dynamics}
\label{p:sec:renewal}

\subsection{The two kernels}

Age-structured population dynamics needs exactly two single-cell
ingredients, and both are already in closed form. The \emph{survival kernel}
is the probability that a cell infected $\alpha$ time units ago is still
productively infected:
\begin{equation}
S(\alpha):=\Ifix(\alpha)=\frac{(a-b)^2\,\ee^{\theta\alpha}}
{\lrb{B+A\ee^{\theta\alpha}}\lrb{a\ee^{\theta\alpha}-b}}.
\label{p:eq:Skernel}
\end{equation}
The \emph{release kernel} is the expected release flux at cell age $\alpha$:
\begin{equation}
g(\alpha):=\delta K(\alpha)=V'(\alpha).
\label{p:eq:gkernel}
\end{equation}
The reasoning behind \cref{p:eq:gkernel} is the one the failed proposal of
\cref{p:sec:obstruction} stated and then did not use: a cell bursts at rate
$\delta X_\alpha$ and releases $X_\alpha$ particles when it does, so the
expected number released in $[\alpha,\alpha+\dd\alpha]$ is
$\delta\,\mean{X_\alpha^2}\,\dd\alpha=\delta K(\alpha)\,\dd\alpha$. With $K$
from \cref{p:eq:K} both kernels are explicit functions of age, which is what
it means for the population model to be \emph{forced by} the intracellular
process (\cref{p:fig:kernels}).

\begin{figure}[tbp]
\centering
\includegraphics[width=0.98\textwidth]{N4b_7_kernels_three_regimes.pdf}
\caption{Survival $\Ifix$ (a)--(c) and release $g=\delta K$ (d)--(f)
across the flooding boundary of \cref{p:tab:trio}. Release rises then
falls --- a cell must accumulate before it bursts --- while survival
decays; shaded areas check $\int_0^\infty\Ifix=\mean{T_{\rm prod}}$ and
$\int_0^\infty g=V_\infty$. Parameters, left to right:
$(\lambda,\mu,\delta)=(1,0,0.1)$, $(1,0.5,1/3)$, $(1,0.9,0.1)$, with
$L=1.10$, $1.00$, $0.42$ and $V_\infty=11.00$, $2.00$, $1.32$. Independent
$g$-scales keep the timing shift visible.}
\label{p:fig:kernels}
\end{figure}

\subsection{The renewal system}

Hold target cells $T$ fixed, let infection occur at rate $\gamma$ per
virion, clear virions at rate $c$, and define the incidence
\begin{equation}
i(t)=\gamma T\,\Vfree(t).
\label{p:eq:incidence}
\end{equation}
Each infection event at time $t-\alpha$ starts a cell of age $\alpha$, still
productive with probability $\Ifix(\alpha)$ and releasing at expected rate
$g(\alpha)$. Summing over infection times and adding the $\Icell_0$ cells
already infected at $t=0$ gives the system.

\begin{definition}[The burst-aware renewal system]
\label{p:def:renewal-system}
\begin{equation}
\Icell(t)=\Icell_0\,\Ifix(t)+\int_0^t i(t-\alpha)\,\Ifix(\alpha)\,\dd\alpha
=\Icell_0\,\Ifix+(i*\Ifix)(t),
\label{p:eq:Icell}
\end{equation}
\begin{equation}
\ddt{\Vfree}(t)
=\Icell_0\,g(t)+\int_0^t i(t-\alpha)\,g(\alpha)\,\dd\alpha-c\,\Vfree(t)
=\Icell_0\,g+(i*g)(t)-c\,\Vfree(t).
\label{p:eq:Vfree}
\end{equation}
\end{definition}

\begin{figure}[t]
\centering
\includegraphics[width=0.98\textwidth]{F4b_1_renewal_schematic.pdf}
\caption{Bookkeeping of \cref{p:def:renewal-system}: infection events start
cohorts; the kernels $\Ifix$ and $g=\delta K$ weight them by age; the
feedback $i=\gamma T\Vfree$ closes the loop.}
\label{p:fig:renewal-schematic}
\end{figure}

\noindent
These are not an approximation to the underlying stochastic process. They
are its first moments exactly, and the reason is that every rate in sight is
linear.

\begin{proposition}[The renewal system is exact in the mean]
\label{p:prop:first-moment}
Consider the stochastic process in which each infection event starts an
independent birth--death--catastrophe cell, each released particle
independently infects at rate $\gamma T$ or clears at rate $c$, and the
$\Icell_0$ cells present at $t=0$ have age zero. Write $\Icell(t)$ and
$\Vfree(t)$ for the expected numbers of productive cells and free particles.
Then $\Icell$ and $\Vfree$ satisfy \cref{p:eq:Icell,p:eq:Vfree} exactly.
\end{proposition}

\begin{proof}
Let $N$ be the point process of infection events. Since each free particle
infects at rate $\gamma T$ and particles are independent, $N$ has
compensator $\gamma T\,\Vfree^{\omega}(t)\,\dd t$, where $\Vfree^\omega$ is
the realised particle count; taking expectations, its mean intensity is
$i(t)=\gamma T\,\Vfree(t)$. A cell founded at time $t-\alpha$ evolves
independently of $N$ and of every other cell, so it is still productive at
$t$ with probability $\Ifix(\alpha)$ and its expected release flux at $t$ is
$g(\alpha)$, both by definition of the two kernels. Conditioning on $N$ and
applying Campbell's formula --- legitimate because $\Ifix$ and $g$ are
bounded and non-negative, so Fubini applies --- gives
$\Icell(t)=\Icell_0\Ifix(t)+(i*\Ifix)(t)$ and an expected release flux of
$\Icell_0 g(t)+(i*g)(t)$. Clearance removes particles at rate $c$ per
particle, contributing $-c\,\Vfree(t)$ to the mean. Collecting the terms
gives \cref{p:eq:Vfree}.
\end{proof}

\noindent
The step that would fail for a nonlinear incidence is the identification of
the mean intensity with $\gamma T\,\Vfree$: for $f(\Vfree,T)$ nonlinear,
$\mean{f(\Vfree^\omega)}\neq f(\mean{\Vfree^\omega})$ and the moment
hierarchy does not close. Assumption (A1), which holds $T$ fixed, is what
keeps the incidence linear. Test~H of \cref{p:app:verification} checks
\cref{p:prop:first-moment} against an independent Gillespie ensemble.

\Cref{p:tab:correspondence} sets the two systems side by side. Both
constants are replaced, not only the release: memoryless removal
$d_{\Icell}\Icell$ becomes the age structure implicit in \cref{p:eq:Icell},
and constant release $p\Icell$ becomes the convolution with $g$. A
replacement touching only the release term --- the first thought of
\cref{p:sec:obstruction} --- would leave the removal side memoryless, and
still wrong.

\begin{table}[htbp]
\centering
\renewcommand{\arraystretch}{1.25}
\small
\begin{tabular}{@{}p{0.46\textwidth}p{0.48\textwidth}@{}}
\toprule
Classical BMVR \cref{p:eq:bmvr-classical}
& Renewal system \crefrange{p:eq:Icell}{p:eq:Vfree} \\
\midrule
$\Icell'=i-d_{\Icell}\Icell$
& $\Icell=\Icell_0\Ifix+i*\Ifix$ \\
constant removal rate $d_{\Icell}$
& age-structured survival via $\Ifix$ \\[0.5em]
$\Vfree'=p\Icell-c\Vfree$
& $\Vfree'=\Icell_0 g+i*g-c\Vfree$ \\
constant release rate $p$
& age-structured release via $g=\delta K$ \\
\bottomrule
\end{tabular}
\caption{What the renewal system replaces. Both constants go, not only the
release rate.}
\label{p:tab:correspondence}
\end{table}

Differentiating \cref{p:eq:Icell} makes the structure visible:
\begin{equation}
\ddt{\Icell}
= i(t)+\Icell_0\,\Ifix'(t)+\int_0^t i(t-\alpha)\,\Ifix'(\alpha)\,\dd\alpha,
\label{p:eq:Icell-diff}
\end{equation}
in which incidence appears as a source exactly as in the classical model,
while the remaining terms encode age-dependent loss of productivity. There
is in general no constant $d_{\Icell}$ reducing this to
$\Icell'=i-d_{\Icell}\Icell$ for every history of $i$: that holds only when
everything grows exponentially, and the exception is what
\cref{p:sec:projection} exploits.

\begin{remark}[What is and is not new]
\label{p:rem:novelty}
The renewal framework is standard --- it is the McKendrick--von Foerster
machinery of mathematical demography
\cite{mckendrick1926applications,vonfoerster1959some,metz1986dynamics} ---
and so are age-structured within-host models: Nelson et al.\
\cite{nelson2004agestructured} write down a within-host system with
age-dependent production and death rates that is structurally the same as
\crefrange{p:eq:Icell}{p:eq:Vfree}, and take the age dependence as given.
What is new here is where the kernels come from. In
\cref{p:eq:Skernel,p:eq:gkernel} the survival and release schedules are
\emph{derived in closed form from a mechanistic intracellular stochastic
process}, so that the population model has no free age structure to fit, and
the consequences drawn in the next two sections --- the effective-parameter
map, the identifiability obstruction, the flooding criterion --- follow from
the intracellular rates alone.
\end{remark}

\subsection{Generality: any pair of kernels will do}
\label{p:sec:generality}

Nothing in \cref{p:eq:Icell,p:eq:Vfree} used the birth--death--catastrophe
process beyond the two kernels it supplied. With incidence
$i(t)=\gamma T\,\Vfree(t)$, or a nonlinear $f(\Vfree,T)$, \emph{any} pair
$(S,g)$ gives
\begin{equation}
\Icell(t)=\Icell_0\,S(t)+(i*S)(t),
\qquad
\ddt{\Vfree}=\Icell_0\,g(t)+(i*g)(t)-c\,\Vfree(t),
\label{p:eq:skeleton}
\end{equation}
with effective parameters $p_{\rm eff}(r)=\Lap{g}(r)/\Lap{S}(r)$ and, for
linear incidence, $R_0=(\gamma T/c)\int_0^\infty g$. This is the chapter's
most general statement, and what makes release mechanisms commensurable:
budding, bursting, multiply infected cells, recurrent partial dumps and
multi-stage eclipse differ only in their $(S,g)$ (\cref{p:tab:sg}). The
first row is the structural check of \cref{p:fig:D-exp}, that exponential
kernels reduce \cref{p:eq:skeleton} exactly to classical BMVR;
\cref{p:sec:reach} works through the rest.

\begin{table}[htbp]
\centering
\renewcommand{\arraystretch}{1.1}
\small
\begin{tabular}{@{}lll@{}}
\toprule
Model & $S(\alpha)$ & $g(\alpha)$ \\
\midrule
Classical BMVR & $\ee^{-d_{\Icell}\alpha}$ & $p\,\ee^{-d_{\Icell}\alpha}$ \\
Absorbing BDC (this chapter) & $\Ifix(\alpha)$ & $\delta K(\alpha)$ \\
MOI-$k$ BDC (\ChDist) & $I(\alpha)^k-D(\alpha)^k$ & $\delta K_k(\alpha)$ \\
Reset catastrophe (\cref{p:sec:reset}) & $\ee^{-d_{\Icell}\alpha}$ & $\delta \ee^{-d_{\Icell}\alpha}\mean{X_\alpha^2}$ \\
HIV multi-stage (\cref{p:sec:hiv}) & stage survival & $p\times$ productive occupancy \\
Two-type / GATE (\ChTwoType) & two-type $\Ifix$ & two-type second-moment flux \\
\bottomrule
\end{tabular}
\caption{One skeleton, many mechanisms. Each row is a different single-cell
story; all of them induce \cref{p:eq:skeleton}, and differ only in the pair
of kernels they supply.}
\label{p:tab:sg}
\end{table}

\subsection{Numerical verification}
\label{p:sec:verification}

The renewal system was checked end to end by a numerical suite at four
levels: kernel integrity, the reduction at $\gamma T=0$, the structural
limits available in closed form, and stochastic ground truth. All 62 checks
pass. Exponential kernels reproduce classical BMVR to relative error below
$5\times10^{-5}$; the late growth rate equals the characteristic root of
\cref{p:sec:projection} to about $10^{-7}$; the $R_0$ threshold is sharp;
and an independent Gillespie ensemble reproduces the renewal means for
$\Icell$ and $\Vfree$ to within $1$--$2\%$ (\cref{p:fig:gillespie}). The
catalogue is \cref{p:app:verification}.

\begin{figure}[tbp]
\centering
\includegraphics[width=0.98\textwidth]{H_gillespie.pdf}
\caption{Stochastic ground truth. Ensemble means of an independent
Gillespie simulation --- birth--death--catastrophe cells with impulse
releases at burst times, exponential clearance, no reinfection --- against
\crefrange{p:eq:Icell}{p:eq:Vfree}. Left column $\mu=0$, right column
$\mu>0$; (a)--(b) infected cells, (c)--(d) free particles. The curves are
indistinguishable at plotting resolution and the relative $L^2$ discrepancy
of $1$--$2\%$ is sampling error: the renewal equations are not an
approximation to the simulation, and this is the check that says so.}
\label{p:fig:gillespie}
\end{figure}

\FloatBarrier
